\section{Introducción}

Este documento presenta el trabajo realizado sobre el caso de estudio 1 de la materia de Computación de Alto Rendimiento (HPC), que consistió en tomar un programa secuencial de multiplicación de matrices e implementar versiones paralelas usando concurrencia y paralelismo a bajo nivel, sin frameworks de alto nivel.

El proyecto se desarrolló en varias etapas. Primero se construyó un \textbf{programa base secuencial} con su propia metodología de medición de tiempo (la versión inicial). Luego se implementaron dos versiones paralelas del mismo programa:

\begin{itemize}
    \item Una versión con \textbf{POSIX Threads (pthreads)}, que reparte el trabajo entre varios hilos dentro de un mismo proceso.
    \item Una versión con \textbf{procesos POSIX (\texttt{fork()})} que reparte el trabajo entre procesos hijos, usando \textbf{memoria compartida POSIX} para que todos vean las mismas matrices.
\end{itemize}

Ambas versiones se compararon con la versión secuencial de referencia usando la métrica estándar de \textit{speedup}. El experimento se diseñó cuidadosamente: se variaron dos parámetros (el tamaño de la matriz $N$ y la cantidad de hilos/procesos $T$), se repitió cada configuración 10 veces para tener resultados confiables, y se midió solo el trabajo de multiplicación (excluyendo reserva de memoria, generación de datos aleatorios e impresión de resultados).

El resto del documento está organizado así: la sección 2 enumera los objetivos; la sección 3 presenta los conceptos teóricos mínimos; la sección 4 describe el programa base secuencial; la sección 5 explica la metodología de medición de tiempo; las secciones 6 y 7 describen las versiones paralelas con hilos y procesos respectivamente; la sección 8 explica el diseño experimental; la sección 9 muestra los resultados; la sección 10 los analiza; y la sección 11 cierra con las conclusiones.